\documentclass[11pt,a4paper,oneside]{report}

\usepackage[hangul]{kotex}
\setmainhangulfont{Malgun Gothic}
\setsanshangulfont{Malgun Gothic}
\setmonohangulfont{Malgun Gothic}

\usepackage[a4paper,margin=2.35cm,top=2.45cm,bottom=2.45cm]{geometry}
\usepackage{setspace}
\linespread{1.12}
\setlength{\parindent}{0pt}
\setlength{\parskip}{0.35em}

\usepackage{amsmath,amssymb}
\usepackage{booktabs,tabularx,array}
\usepackage{enumitem}
\usepackage{xcolor}
\usepackage{tikz}
\usetikzlibrary{arrows.meta,positioning}
\usepackage[most]{tcolorbox}
\usepackage{listings}
\usepackage{hyperref}
\usepackage{url}

\definecolor{floorGreen}{RGB}{70,180,90}
\definecolor{leftBlue}{RGB}{60,150,255}
\definecolor{rightOrange}{RGB}{255,160,60}
\definecolor{frontPink}{RGB}{230,80,160}
\definecolor{otherGray}{RGB}{150,150,150}
\definecolor{codebg}{RGB}{248,248,245}
\definecolor{coderule}{RGB}{210,210,210}

\newtcolorbox{keybox}[1][]{colback=blue!4,colframe=blue!45!black,sharp corners,boxrule=0.5pt,#1}
\newtcolorbox{warnbox}[1][]{colback=orange!6,colframe=orange!55!black,sharp corners,boxrule=0.5pt,#1}

\lstdefinestyle{base}{
  basicstyle=\ttfamily\footnotesize,
  backgroundcolor=\color{codebg},
  frame=single,
  rulecolor=\color{coderule},
  showstringspaces=false,
  breaklines=true,
  columns=fullflexible,
  keepspaces=true,
  tabsize=2
}
\lstset{style=base}

\title{\textbf{KITTI LiDAR Floor/Wall Segmentation 원리 정리}\\[0.35em]
\large Velodyne point cloud 기반 바닥, 좌/우/정면 벽 후보 분류}
\author{LC-EKF VIO project notes}
\date{2026-05-04}

\begin{document}
\maketitle
\tableofcontents

\chapter{왜 LiDAR로 나누었나}

이번 구현의 목적은 KITTI raw 데이터에서 바닥, 좌측 벽, 우측 벽, 정면 벽 후보를 Rerun에서 색으로 확인하는 것이다.

기존 VIO frontend는 stereo image에서 sparse feature를 추적한다. 이 feature map은 pose 추정에는 충분하지만, 면을 안정적으로 분할하기에는 점 밀도가 낮고 표면 전체를 덮지 않는다. 반면 KITTI raw에는 각 camera frame과 동기화된 Velodyne point cloud가 있으므로, 3D geometry를 직접 사용하면 바닥과 벽 후보를 더 단순하고 안정적으로 분리할 수 있다.

\begin{keybox}
현재 구현은 semantic segmentation이 아니다. 딥러닝 모델 없이 3D point cloud의 위치, 평면성, 법선 방향만 이용하는 geometry 기반 plane segmentation이다.
\end{keybox}

\chapter{입력 데이터와 좌표계}

KITTI Velodyne \path{.bin} 파일은 점 하나당 네 개의 \texttt{float}로 저장된다.

\begin{lstlisting}
x y z intensity
\end{lstlisting}

Velodyne 좌표계는 다음 기준으로 사용했다.

\begin{center}
\begin{tabularx}{0.8\linewidth}{@{}cX@{}}
\toprule
축 & 의미 \\
\midrule
$x$ & 차량 전방 \\
$y$ & 차량 좌측 \\
$z$ & 위쪽 \\
\bottomrule
\end{tabularx}
\end{center}

따라서 바닥은 대체로 낮은 $z$에 있고 법선이 $z$축에 가깝다. 좌/우 벽은 수직면이므로 법선의 $z$ 성분이 작고, 좌우 방향 벽은 법선이 주로 $y$축 방향이다. 정면 벽은 법선이 주로 $x$축 방향이다.

\chapter{전체 처리 파이프라인}

\begin{center}
\begin{tikzpicture}[
  node distance=0.9cm,
  box/.style={draw, rounded corners=2pt, align=center, minimum width=4.1cm, minimum height=0.8cm, fill=blue!4},
  arrow/.style={-{Latex[length=2.2mm]}, thick}
]
\node[box] (load) {Velodyne \path{.bin} 로드\\\texttt{x,y,z,intensity}};
\node[box, below=of load] (roi) {ROI 필터링\\거리/측방/높이 범위 제한};
\node[box, below=of roi] (floor) {바닥 RANSAC\\낮은 $z$, $n_z$ 큰 plane};
\node[box, below=of floor] (walls) {남은 점에서 wall RANSAC\\left/right/front 후보군 분리};
\node[box, below=of walls] (labels) {라벨 및 색상 부여\\Rerun Points3D 또는 PLY 저장};

\draw[arrow] (load) -- (roi);
\draw[arrow] (roi) -- (floor);
\draw[arrow] (floor) -- (walls);
\draw[arrow] (walls) -- (labels);
\end{tikzpicture}
\end{center}

\section{ROI 필터}

먼저 너무 멀거나 의미 없는 높이의 점을 제외한다. 기본 범위는 다음과 같다.

\begin{center}
\begin{tabularx}{0.9\linewidth}{@{}c c X@{}}
\toprule
항목 & 기본값 & 의미 \\
\midrule
$x$ & $[-5, 50]$ m & 차량 뒤쪽 일부와 전방 50 m까지 사용 \\
$y$ & $[-25, 25]$ m & 좌우 25 m 범위 사용 \\
$z$ & $[-3, 3]$ m & 너무 낮거나 높은 점 제거 \\
\bottomrule
\end{tabularx}
\end{center}

\section{평면 모델}

모든 plane fitting은 다음 평면식을 사용한다.

\[
\mathbf{n}^{\mathsf{T}}\mathbf{p} + d = 0,
\qquad
\|\mathbf{n}\| = 1
\]

세 점 $\mathbf{p}_a,\mathbf{p}_b,\mathbf{p}_c$를 뽑으면 법선은 cross product로 계산한다.

\[
\mathbf{n}
= \frac{(\mathbf{p}_b-\mathbf{p}_a) \times (\mathbf{p}_c-\mathbf{p}_a)}
       {\|(\mathbf{p}_b-\mathbf{p}_a) \times (\mathbf{p}_c-\mathbf{p}_a)\|}
\]

점 $\mathbf{p}$가 평면에 속하는지는 거리로 판단한다.

\[
\operatorname{dist}(\mathbf{p}, \Pi)
= |\mathbf{n}^{\mathsf{T}}\mathbf{p} + d|
\]

거리 threshold 안에 들어오는 점을 inlier로 센다.

\section{RANSAC}

RANSAC은 다음 반복을 수행한다.

\begin{enumerate}[leftmargin=2em]
\item 후보 점에서 세 점을 무작위로 선택한다.
\item 세 점이 만드는 평면을 계산한다.
\item 그 평면의 법선 방향이 원하는 class 조건과 맞는지 검사한다.
\item 모든 후보 점에 대해 평면 거리 threshold 이내인 inlier를 센다.
\item 가장 많은 inlier를 가진 평면을 채택한다.
\end{enumerate}

구현에서는 deterministic하게 같은 결과가 나오도록 고정 seed를 사용했다.

\chapter{라벨별 판단 기준}

\section{Floor}

바닥은 낮은 $z$ 범위에서 후보를 모은 뒤, 법선이 위쪽에 가까운 평면을 찾는다.

\[
|n_z| \ge 0.85
\]

또한 원점 근처의 바닥 높이

\[
z_0 = -\frac{d}{n_z}
\]

가 floor search 범위 안에 있어야 한다.

\section{Left/Right Wall}

좌/우 벽 후보는 floor로 라벨링되지 않은 점 중 수직 구조일 가능성이 있는 점에서 찾는다.

\begin{itemize}[leftmargin=2em]
\item left wall 후보: $y \ge 2.0$ m
\item right wall 후보: $y \le -2.0$ m
\item vertical plane 조건: $|n_z| \le 0.25$
\item side wall 방향 조건: $|n_y| \ge 0.70$
\end{itemize}

추가로 너무 작은 cluster가 벽으로 잡히지 않도록, inlier 수와 horizontal/vertical extent 조건을 검사한다.

\section{Front Wall}

정면 벽은 전방 영역의 수직면을 찾는다.

\begin{itemize}[leftmargin=2em]
\item 후보: $x \ge 5.0$ m, $|y| \le 12.0$ m
\item vertical plane 조건: $|n_z| \le 0.25$
\item front wall 방향 조건: $|n_x| \ge 0.70$
\end{itemize}

\begin{warnbox}
일반 도로 KITTI 장면에서는 정면에 큰 수직 벽이 항상 존재하지 않는다. 따라서 \texttt{front\_wall}은 장면에 따라 적거나, 차량/가드레일/수풀 일부가 섞일 수 있다.
\end{warnbox}

\chapter{색상과 Rerun 출력}

라벨별 색상은 \path{segment_color_rgb()}에 정의했다.

\begin{center}
\begin{tabularx}{0.95\linewidth}{@{}l c c X@{}}
\toprule
Label & Color & RGB & 의미 \\
\midrule
\texttt{floor} & \textcolor{floorGreen}{green} & (70, 180, 90) & 바닥 plane inlier \\
\texttt{left\_wall} & \textcolor{leftBlue}{blue} & (60, 150, 255) & 차량 좌측 벽 후보 \\
\texttt{right\_wall} & \textcolor{rightOrange}{orange} & (255, 160, 60) & 차량 우측 벽 후보 \\
\texttt{front\_wall} & \textcolor{frontPink}{pink} & (230, 80, 160) & 차량 전방 벽 후보 \\
\texttt{other} & \textcolor{otherGray}{gray} & (150, 150, 150) & 나머지 점 \\
\bottomrule
\end{tabularx}
\end{center}

Rerun에서는 하나의 colored \texttt{Points3D}로 기록한다.

\begin{lstlisting}
/lidar/segments
/metrics/lidar_floor_points
/metrics/lidar_left_wall_points
/metrics/lidar_right_wall_points
/metrics/lidar_front_wall_points
\end{lstlisting}

Rerun을 쓰지 않는 경우에는 같은 색상 정보를 가진 PLY를 저장한다.

\begin{lstlisting}
<output>/lidar_segments/*.ply
\end{lstlisting}

\chapter{코드 위치}

\begin{center}
\begin{tabularx}{\linewidth}{@{}lX@{}}
\toprule
파일/함수 & 역할 \\
\midrule
\path{include/lidar/lidar_segmenter.hpp} & point, label, option, result 타입 정의 \\
\path{src/lidar/lidar_segmenter.cpp} & KITTI \path{.bin} 로드, RANSAC segmentation, PLY 저장, 색상 매핑 \\
\path{LidarSegmenter::process()} & floor, left/right/front wall 라벨링 핵심 로직 \\
\path{segment_color_rgb()} & Rerun/PLY에 들어가는 라벨별 색상 \\
\path{src/io/kitti_raw_reader.cpp} & \path{velodyne_points/data/*.bin} path 수집 \\
\path{apps/run_vio.cpp} & CLI 옵션, main loop 호출, Rerun logging \\
\path{RerunLogger::log_lidar_segments()} & \path{/lidar/segments} Points3D 기록 \\
\bottomrule
\end{tabularx}
\end{center}

\chapter{실행 명령}

Rerun-enabled build:

\begin{lstlisting}
wsl bash -lc 'cd /mnt/d/02_research/04_cpp_LC-EKF_VIO && cmake -S . -B build_verify_rerun -DCMAKE_BUILD_TYPE=Release -DENABLE_RERUN=ON && cmake --build build_verify_rerun -j$(nproc)'
\end{lstlisting}

전체 frame 실행:

\begin{lstlisting}
wsl bash -lc 'cd /mnt/d/02_research/04_cpp_LC-EKF_VIO && mkdir -p results/kitti_raw_2011_09_26_drive_0117_full && ./build_verify_rerun/run_vio config/kitti_raw_2011_09_26_drive_0117_local_max_frame.yaml --segment-lidar --segment-every 1 --rerun-save results/kitti_raw_2011_09_26_drive_0117_full/lidar_segments_full.rrd --rerun-image-every 10'
\end{lstlisting}

Rerun 열기:

\begin{lstlisting}
rerun D:\02_research\04_cpp_LC-EKF_VIO\results\kitti_raw_2011_09_26_drive_0117_full\lidar_segments_full.rrd
\end{lstlisting}

\chapter{직접 수집한 데이터로 응용하기}

현재 구현은 \texttt{run\_vio}의 KITTI raw reader에 연결되어 있다. 따라서 가장 빠른 응용 방법은 직접 수집한 데이터를 KITTI raw와 비슷한 폴더 구조로 변환하는 것이다.

\section{수집해야 하는 데이터}

\begin{center}
\begin{tabularx}{\linewidth}{@{}lX@{}}
\toprule
데이터 & 필요 이유 \\
\midrule
Stereo camera image & 현재 VIO frontend가 left/right image를 필요로 한다. \\
IMU 또는 OXTS-like packet & 현재 reader가 IMU 초기화와 propagation 입력을 필요로 한다. \\
LiDAR point cloud & floor/wall segmentation의 실제 입력이다. \\
Calibration YAML & camera intrinsics, stereo extrinsics, camera--IMU transform을 설정한다. \\
\bottomrule
\end{tabularx}
\end{center}

\begin{warnbox}
LiDAR segmentation 자체는 LiDAR 점만 있으면 가능하지만, 현재 실행 파일은 camera/IMU도 함께 로드한다. LiDAR만 단독으로 돌리고 싶다면 \path{LidarSegmenter}를 호출하는 별도 LiDAR-only executable을 추가하는 것이 더 깔끔하다.
\end{warnbox}

\section{권장 폴더 구조}

직접 수집한 데이터를 다음처럼 정리한다. 파일명은 사전순 정렬이 곧 frame 순서가 되도록 10자리 zero padding을 권장한다.

\begin{lstlisting}
data/my_dataset/my_drive_sync/
  image_00/
    timestamps.txt
    data/
      0000000000.png
      0000000001.png
  image_01/
    timestamps.txt
    data/
      0000000000.png
      0000000001.png
  oxts/
    timestamps.txt
    data/
      0000000000.txt
      0000000001.txt
  velodyne_points/
    data/
      0000000000.bin
      0000000001.bin
\end{lstlisting}

현재 \path{KittiRawReader}는 Velodyne timestamp 파일은 쓰지 않고, camera frame index와 같은 index의 \path{.bin} 파일을 사용한다. 즉, \texttt{image\_00/data/0000000042.png}에 대응하는 LiDAR는 \texttt{velodyne\_points/data/0000000042.bin}이라고 가정한다.

\section{LiDAR bin 변환 규칙}

각 LiDAR frame은 little-endian float32 네 개를 순서대로 저장한다.

\begin{lstlisting}
float32 x
float32 y
float32 z
float32 intensity
\end{lstlisting}

좌표계는 KITTI Velodyne 기준으로 맞춘다.

\begin{center}
\begin{tabularx}{0.85\linewidth}{@{}cX@{}}
\toprule
축 & 직접 데이터 변환 규칙 \\
\midrule
$x$ & 차량/센서 전방이 양수 \\
$y$ & 차량/센서 좌측이 양수 \\
$z$ & 위쪽이 양수 \\
단위 & meter \\
\bottomrule
\end{tabularx}
\end{center}

예를 들어 사용하는 LiDAR SDK가 $x$ right, $y$ forward, $z$ up 같은 좌표계를 낸다면, \path{.bin} 저장 전에 KITTI 좌표계로 회전/축 교환을 해야 한다.

\section{Timestamp와 동기화}

현재 reader는 timestamp line에서 시각 부분을 읽는다.

\begin{lstlisting}
2011-09-26 15:38:03.123456789
\end{lstlisting}

직접 데이터에서는 절대 날짜보다 frame 순서와 sensor 간 동기화가 중요하다.

\begin{itemize}[leftmargin=2em]
\item camera left/right frame 수가 같아야 한다.
\item Velodyne \path{.bin} frame 수가 camera frame 수와 같거나 더 많아야 한다.
\item 현재 구현은 nearest timestamp matching이 아니라 같은 index matching을 사용한다.
\item frame drop이 있으면 image와 LiDAR 파일명을 같은 index로 다시 정렬한다.
\end{itemize}

\section{Config 작성}

기존 config를 복사해서 dataset/output/calibration을 바꾼다.

\begin{lstlisting}
Copy-Item `
  D:\02_research\04_cpp_LC-EKF_VIO\config\kitti_raw_2011_09_26_drive_0117_local_max_frame.yaml `
  D:\02_research\04_cpp_LC-EKF_VIO\config\my_drive.yaml
\end{lstlisting}

수정할 항목:

\begin{lstlisting}
dataset_type: "kitti_raw"
dataset: "/mnt/d/02_research/04_cpp_LC-EKF_VIO/data/my_dataset/my_drive_sync"
output:  "/mnt/d/02_research/04_cpp_LC-EKF_VIO/results/my_drive"
max_frames: 0
\end{lstlisting}

그리고 \texttt{cam0}, \texttt{cam1}의 intrinsics와 \path{T_cam_imu}를 직접 장비 calibration으로 바꾼다.

\section{실행}

\begin{lstlisting}
wsl bash -lc 'cd /mnt/d/02_research/04_cpp_LC-EKF_VIO && mkdir -p results/my_drive && ./build_verify_rerun/run_vio config/my_drive.yaml --segment-lidar --segment-every 1 --rerun-save results/my_drive/lidar_segments.rrd --rerun-image-every 10'
\end{lstlisting}

Rerun으로 확인한다.

\begin{lstlisting}
rerun D:\02_research\04_cpp_LC-EKF_VIO\results\my_drive\lidar_segments.rrd
\end{lstlisting}

\section{현장 수집 체크리스트}

\begin{itemize}[leftmargin=2em]
\item 차량/센서가 멈춰 있을 때 몇 초간 기록해서 IMU 초기화가 안정되게 한다.
\item LiDAR가 바닥을 충분히 보도록 설치한다. 너무 하늘/전방만 보면 floor RANSAC이 약해진다.
\item 좌우 벽을 보고 싶다면 도로 가장자리, 복도, 주차장 벽처럼 수직 구조가 있는 장면을 포함한다.
\item frame drop이 생기면 camera/LiDAR index alignment를 먼저 고친다.
\item Rerun에서 `/lidar/segments`가 기울어져 보이면 LiDAR 좌표계 변환이 잘못된 것이다.
\end{itemize}

\chapter{현재 한계}

\begin{itemize}[leftmargin=2em]
\item geometry 기반이므로 semantic class를 이해하지 않는다.
\item 벽과 비슷한 수직 구조물(차량 측면, 가드레일, 수풀)이 wall 후보로 섞일 수 있다.
\item 각 frame을 독립적으로 처리하므로 temporal smoothing은 없다.
\item 현재 segmentation 결과는 EKF update에 사용하지 않고 시각화/분석용으로만 사용한다.
\item 현재 실행 파일은 KITTI-like stereo/IMU/LiDAR 구조를 가정한다. 임의 LiDAR-only 로그를 바로 읽지는 않는다.
\end{itemize}

\end{document}
